La sucesión de Fibonacci se define como:

\begin{equation}
    f_n = f_{n-1} + f_{n-2}
\end{equation}


\noindent donde $f_0 = 0$ y $f_1 = 1$. Por lo tanto, para el siguiente valor de $f$ es:

\begin{equation*}
    f_2 = 0 + 1 = 1
\end{equation*}

\noindent Para $f_3$

\begin{equation*}
    f_3 = f_2 + f_1 = 1 + 1 = 2
\end{equation*}

\noindent y así sucesivamente.

El siguiente código de MATLAB representa la función de Fibonacci para cualquier $n$.

\begin{figure}[ht!]
    \centering
    \begin{lstlisting}[caption = {La función \texttt{fibonacci2} calcula el valor de la sucesión para un valor $n$}, label = {cod:1}]
    function fn = fibonacci2(n)
        a = 0; b = 1;
        arr = 1:1:n;
        for i = arr
            fn = a + b;
            temp = a;
            a = fn ; b = temp;
        end
    end
    \end{lstlisting}
\end{figure}

Si llamamos la función y le ponemos el valor de $n = 40$, el valor obtenido es $102334155$. Los siguientes códigos representan lo mismo pero usando bucles \texttt{while}.


\begin{lstlisting}[caption = {La función \texttt{fibonacci} representa la función para calcular los primeros $n$ valores de la sucesión. \texttt{fibonacci\_10000} calcula el valor de la sucesión una vez que es mayor a 10000. Por último, \texttt{fibonacci\_v} le pasamos como argumento un valor $v$ donde calculará el último valor de la sucesión cuando $f_n > v$.}, label = {cod:2}]
function fn = fibonacci(n)
    a = 0 ; b = 1; %f_0 = 0,  f_1 = 1
    m = n - 1; %le resto 1 ya que arranco con 2 valores de la sucesión
    i = 0;
    while i <= m
        fn = a + b;
        i = i + 1;
        temp = a;
        a = fn; b = temp;
    end
end


function fn = fibonacci_10000()
    a = 0 ; b = 1;
    fn = a + b;
    while fn < 10000
        fn = a + b;
        temp = a;
        a = fn ; b = temp;
    end
end


function fn = fibonacci_v(v)
    a = 0; b = 1;
    i = 1;
    fn = 0;
    while fn < v
        fn = a + b;
        i = i + 1;
        temp = a;
        a = fn ; b = temp;
    end
end
\end{lstlisting}



\begin{tcolorbox}[colframe = white!40!black, colback = black!10!white, sharp corners, breakable, title = Criterio de estabilidad de Routh-Hurwitz]

El criterio de estabilidad de Routh-Hurwitz dice si existen o no raíces inestables en una ecuación polinomial, sin tener que obtenerlas en realidad. Este criterio de estabilidad sólo se aplica a los polinomios con una cantidad finita de términos. 

El procedimiento en el criterio de estabilidad de Routh es el siguiente:

\begin{enumerate}
    \item Se escribe el denominador de la función transferencia de la forma siguiente:
    
    \begin{equation}
        a_0s^n + a_1s^{n-1} + \cdots + a_{n-1}s + a_n = 0
    \end{equation}

    \noindent donde los coeficientes son cantidad reales. Se supone que $a_n \neq 0$; es decir, se elimina cualquier raíz cero.

    \item Si todos los coeficientes son positivos, se ordenan los coeficientes del polinomio en filas y columnas con el patrón siguiente:
    
    \begin{align*}
        s^n  &\qquad a_0 \quad a_2 \quad a_4 \quad a_6 \cdots \\
        s^{n-1}  &\qquad a_1 \quad a_3 \quad a_5 \quad a_7 \cdots \\
        s^{n-2}  &\qquad b_1 \quad b_2 \quad b_3 \quad b_4 \cdots \\
        s^{n-3}  &\qquad c_1 \quad c_2 \quad c_3 \quad c_4 \cdots \\
        \cdot &\qquad \cdot\\
        \cdot &\qquad \cdot\\
        \cdot &\qquad \cdot\\
        s^2  &\qquad e_1 \quad e_2 \\
        s^1  &\qquad f_1\\
        s^0  &\qquad g_1\\
    \end{align*}

    \noindent El proceso de armar los valores de $b$, $c \cdots$, $e$, $f$ y $g$ se realiza mediante un proceso parecido a tomar el determinante de una matriz $2\times2$. 
\end{enumerate}


El criterio asegura que si existe algún cambio de signo en la primera columna de los coeficientes, entonces el sistema es inestable (tiene una raíz en el semiplano derecho).

\end{tcolorbox}

El siguiente ejemplo muestra cómo aplicar el criterio para un sistema de 4to orden:

\begin{equation}
    y^{(4)} + 2\dddot{y} + 4\ddot{y} + \dot{y} + 1 = r(t)
\end{equation}

La matriz $H$ es el arreglo de coeficientes definido anteriormente:

\begin{equation*}
    H =
    \begin{bmatrix}
        1 & 4 & 1 \\
        2 & 1 & 0 \\
        \alpha_1 & \alpha_2 & 0 \\
        \beta_1 & 0 & 0 \\
        \gamma_1 & 0 & 0 \\
    \end{bmatrix}
\end{equation*}

\noindent donde 

\begin{equation*}
    \alpha_1 = 3.5, \qquad \alpha_2 = 1, \qquad \beta_1 = 0.4286, \qquad \beta_2 = 0, \qquad \gamma_1 = 3.5
\end{equation*}


El siguiente código representa la función que evalúa si el sistema es o no estable:


\begin{lstlisting}[caption = {Criterio de Routh aplicado al sistema de 4to orden}, label = {cod:3}]
function estable(M)
    first_row = M(1);
    for i = 1:length(first_row)
        if i < 0
            disp('Sistema no estable!')
            break
        end
    disp('Sistema estable!')
    end
end

function RH_4(y)
    len = length(y);
    H = zeros(len, 3);
    a4 = y(1); a3 = y(2); a2 = y(3); a1 = y(4); a0 = y(5);
    alpha1 = ((a3*a2) - (a4*a1))/a3;
    alpha2 = ((a3*a0) - (a4*0))/a3;
    beta1 = ((alpha1*a1)-(a3*alpha2))/alpha1;
    gamma1 = ((beta1*alpha1)-(0*alpha1))/beta1;

    H(1, [1:3]) = [a4, a2, a0];
    H(2, [1:2]) = [a3, a1];
    H(3, [1:2]) = [alpha1, alpha2];
    H(4, 1) = beta1;
    H(5, 1) = gamma1;
    
    estable(H);
    H
end

coef1 = [1, 2, 4, 1, 1];
% coef2 = [3, 2, 2, 0, -1];
RH_4(coef1)

\end{lstlisting}



Por último, veamos el siguiente problema. Si $p$ es un número primo se puede obtener un nuevo número primo $q$ de la siguiente manera:

\begin{equation}
    q = 2\cdot3\cdot5\cdot7\cdots p + 1 
\end{equation}


El siguiente código implementa un bucle para buscar el último número primo que puede manejar MATLAB, el cual es: $2.811\times10^{307}$.

\begin{lstlisting}[caption = {}, label = {}]
for i = 300:900
    q = prod(primes(i)) + 1;
    if ~isfinite(q)
            q = prod(primes(i - 1) + 1);
            break
    end
end
\end{lstlisting}







